%!TEX program = lualatex
\documentclass[ngerman,headheight=70pt]{scrartcl}
\usepackage[ngerman]{babel}
\usepackage{blindtext}
\usepackage[utf8]{luainputenc}
\usepackage{scrlayer-scrpage}
\usepackage[T1]{fontenc}
\usepackage{fontspec}
\usepackage{multicol}
\usepackage{enumitem}
\usepackage[german=quotes]{csquotes}
\usepackage[bottom=2cm,footskip=8mm]{geometry}
\setmainfont{Source Sans Pro}
\MakeOuterQuote{"}

\setlength{\parindent}{0pt}
\setlength{\parskip}{6pt}

\pagestyle{plain.scrheadings}
\thispagestyle{scrheadings}
\lohead{\vspace*{-1.5cm}\includegraphics[scale=0.4]{../images/stupa_logo_rot.pdf}}
\rohead{\\Präsidium des\\Studierendenparlaments}

\lofoot[Universität Hamburg · Präsidium des Studierendenparlaments\\\hrule Von-Melle-Park 5 · D-20146 Hamburg]{Universität Hamburg · Präsidium des Studierendenparlaments\\\hrule Von-Melle-Park 5 · D-20146 Hamburg}
\cofoot[]{}
\cohead[\thepage]{}

\renewcommand{\thesubsection}{TOP \arabic{subsection}}
\renewcommand{\thesubsubsection}{\arabic{subsubsection}.}
\setcounter{subsection}{-1}

\begin{document}
    UHH · StuPa-Präsidium · Von-Melle-Park 5 · D-20146 Hamburg

    %%% PAGE ONE (EDIT LATER) %%%%
    % Hier den Titel des Protokolls eintragen
    \section*{Protokoll der 13. Sitzung des Studierendenparlaments vom 12. Januar 2017}

    % Protokollanten eintragen
    \textbf{Protokoll: Jim Martens}\\
    \textbf{Ort: VMP5 0077}\\
    \textbf{Beginn: 20.54 Uhr}\\
    \textbf{Ende: 23.42 Uhr}

    % Wohl die komplizierteste Stelle im Protokoll
    % Als Faustregel sollten nicht mehr als 3 Namen auf einer Zeile stehen
    % Einfach an dieser Vorlage orientieren
    % Im Idealfall ist hier nicht viel zu ändern, da kaum jemand fehlt.
    % Rücktritte müssen natürlich pro Protokoll angepasst werden
    \vspace{0.5cm}
    \begin{tabular}{ll}
        Anwesend: & \\
            RCDS (5 Sitze): & Ramin Shakiba, Jennifer Maack \\
             CampusGrün (14 Sitze): & Laura Franzen, Geoffrey Youett, Elena Rysikova, \\
                                   & Philipp Droll, Yasemin Günther, Melf Johannsen, \\
                                   & Mario Moldenhauer, Christian Lagod, Jim Martens, \\
                                   & Svenja Horn, Damian Hofmann, Armin Günther \\
             BGZM-Liste (2 Sitze) : & \\
             WiWi (2 Sitze): & \\
             Unicorns (5 Sitze): & Andreas Hartkamp, Marielle Hermstrüwer \\
             Liste LINKS (3 Sitze): & Gunhild Berdal, Till Petersen, Sinah Mielich \\
             HWP (2 Sitze): &  \\
             MIN (4 Sitze): &  \\
             SDS* (3 Sitze): & Mena Winkler, Jacob Petersein, Artur Brückmann \\
             Bart-LISTE (2 Sitze): & Timo Zeimet, Malte Peters \\
             LHG (1 Sitz): & Tobias Heisig \\
             harte zeiten (1 Sitz): & Tobias Berking \\
             Jura (1 Sitz): & \\
             AL (2 Sitze): & Karima Schulze, Henri Weber \\
            & \\
        Entschuldigt: & Ramon Weilinger (RCDS), Sonja Lattwesen, Martin Sievert (CG), \\
                      & Jakob Pape (BGZM), Katharina Kucza, Johannes Peplow, \\
                      & Annkathrin Löffler (Unicorns), Ailina Salten, Lotte Rullkötter, \\
                      & Nabil Jabareen, Jan Detampel (MIN) \\
                      &\\
        Unentschuldigt abwesend: & Benjamin Welling (RCDS), Karen Martirosian (BGZM), Claas-Friso Hente,\\
                                & Charleen Lorenz (WiWi), Hakan Tasdemir, Gökce Celik (HWP), \\
                                & Maurice Martens (Jura)\\
        Rücktritte: & Dean Haltmeyer \rightarrow Heiner Rolfs \rightarrow \\
                    & Marie Wollenschläger \rightarrow Damian Hofmann (CG), \\
                    & Ines Friedrich \rightarrow Ronja Winkhardt-Enz \rightarrow \\
                    & Dario Paschke \rightarrow Sonja Lattwesen (CG) \\
        Mandatsverluste: & \\
    \end{tabular}
    \newpage
    %%% Vorgeschlagene Tagesordnung %%%
    \underline{Vorgeschlagene Tagesordnung}
    \begin{enumerate}[label={\textbf{Top \theenumi}},leftmargin=*]
        \item Beitragsordnung (V1617-046) (30 Min.)
        \item VG-Wort (V1617-047) (30 Min.)
        \item Englischsprachige Wahlwerbung (V1617-048) (30 Min.)
        \item Manifest zu G20 (V1617-049) (45 Min.)
        \item Verschiedenes (5 Min.)
    \end{enumerate}

    \newpage

    % Hier beginnt Teil A. Bei Fortsetzungen kommt hier nur der Teil über
    % die Feststellung der endgültigen Fassung der Tagesordnung.
    % Für sonstige Sitzungen ist dies bereits das Gerüst für die Formalia.
    % Man muss nur noch die protokollierten Inhalte eintragen.

    {\Large \underline{Teil A}}

    \subsection{Formalia}

    \textit{
    Gunhild eröffnet die Sitzung um 20.54 Uhr. Ramin vertritt Ramon im Präsidium.
    }

    \subsubsection{Geschäftsbericht Präsidium}

    Gunhild berichtet aus dem Präsidium.

    \begin{itemize}
        \item Beschäftigung mit der Wahlvorbereitung
        \item Wahlhelfer*innen haben Wahlzeitungen und Plakate verteilt - auch
              an Außenstandorte und Wohnheime
        \item Letzte Woche wurden die fristgerechten Briefe von der Briefwahl
              gezählt. 3250 Briefe enthalten potentiell gültige Stimmen. 335 Briefe
              waren eindeutig ungültig. 7,68\% Wahlbeteiligung durch gültige
              Stimmen. Eine starke Beteiligung bei der Urnenwahl ist daher
              wichtig, um die Wahlbeteiligung vom letzten Jahr zu erreichen.
        \item Es sei wichtig die Leute aufzufordern, an der Urne zu überprüfen, ob
              ihre Briefwahlstimme angekommen ist.
        \item Die Urnen werden morgen zusammen mit Herrn Drexler versiegelt. Am
              Montag geht es mit der Urnenwahl los.
        \item Es gibt ein Sonderproblem mit dem Studienkolleg, welches weit weg
              vom Hauptcampus ist. Infomaterial wurde jedoch dort hingebracht
              und am Freitag, 14.1., findet dort ein Infotreff statt.
        \item Gunhild ruft zur Besetzung von Urnen auf. Wenn Leute in FSRen aktiv
              sind, sollte dort für die Wahl mobilisiert werden. Ebenfalls soll
              auch an der Wahl zum AS teilgenommen werden.
        \item Ungefähr 10 AS-Briefumschläge wurden in Briefen bei der Briefzählung
              gefunden. AS-Briefwahlumschläge können an den Urnen und beim Präsidium
              bis Dienstag, 17.1., 13 Uhr abgegeben werden. Anschließend werden
              die AS-Briefe fristgerecht zum Wahlamt der Uni gebracht.
    \end{itemize}

    Jim verliest die Rücktritte. Es gibt einen Vorschlag den Rücktrittsanteil einer
    Liste in die Wahlzeitung aufzunehmen.

    \subsubsection{Anfragen an das Präsidium}

    Philipp merkt an, dass viele Plakate zerstört wurden und fragt, ob auch andere
    Gruppen betroffen sind.

    Karima merkt an, dass auch bei ihnen Plakate kaputt gegangen sind. Sie vermutet
    jedoch eher eine betrunkene Bande.

    Tobias B. hat auch einige Verluste zu beklagen, aber nichts Systematisches.

    Laura merkt an, dass z.B. um das Geomatikum alle Plakate restlos verschwunden
    sind.

    Melf weist auf Orte hin, wo Hausmeister Plakate wegräumen. Um das Geomatikum
    ist das jedoch eher nicht der Fall.

    Paula merkt an, dass Richtung Dammtor und Hauptcampus alles um Silvester
    weg war.

    Karima merkt an, dass man aufpassen sollte, ob Plakate von gewissen Listen
    gezielt entfernt werden.

    Gunhild bietet an, dass der Beschluss "Fairer Wahlkampf" stärker auf der
    Website platziert werden kann. Die Hausmeister könnten wir informieren.

    Franzi weist darauf hin, dass die Hausmeister der HWP über die bestellten
    Tische informiert haben. Gunhild klärt auf, dass die Tische nicht vom
    Präsidium zu den Gebäuden geschleppt werden müssen, sondern direkt zu den
    Gebäuden geliefert werden.

    \subsubsection{Geschäftsbericht AStA}

    \textit{Zeitansage: 21.15 Uhr}

    Franziska beginnt den Geschäftsbericht:

    Ein Dies Academicus wird im WS 17/18 über das Thema Finanzen stattfinden.
    Der AStA soll der Uni sagen, mit welchen Leuten die Vorbereitung zum Dies
    stattfinden soll. Die Vorbereitungsgruppe wollte sich bereits im Januar
    konstituieren.

    Es wurde ein Vertrag mit der Sprinkenhof GmbH und Behörde für Wissenschaft,
    Forschung und Gleichstellung unterzeichnet:
    Der Überseering 35 wird die Interimsunterbringung für den Philturm sein.
    Der Bautermin bzgl. der Philturm-Sanierung steht. Die Sanierung beginnt somit
    im Sommer 2017. Die Bibliotheken und der gesamte Philturm kann in das
    Interimsgebäude ziehen. Es wird aufgrund der großen Entfernung zum Hauptcampus
    eine Herausforderung für studentische und politische Kultur mit sich bringen.

    Die Kantine im Interimsgebäude wird privatwirtschaftlich betrieben. Dem
    Studierendenwerk entfällt die Philmensa. Im Januar wird es Vollversammlung
    des Philturms unter Leitung der Uni geben.

    Die RPO wurde kontrovers im AS diskutiert. "Rechter Pol" musste beteuern,
    dass sie nichts gegen Studienreform hätten. Herr Ritter und Frau Rupp haben
    versprochen einen Entwurf für eine General Studies Umsetzung bis zur RPO-Tagung
    zur Verfügung zu stellen.

    Der AS hat beschlossen, dass in 2. Sitzung der neuen Legislatur erneut eine 1.
    Lesung eines RPO-Entwurfs stattfindet. Den Dies zur Studienreform im Mai gibt
    es weiterhin. Herr Lenzen hat betont, dass eine RPO notwendig ist und dass es sie
    geben wird.

    Armin fährt fort mit dem Bericht:
    Am Dienstag hat die 3. Sitzung der G20-Plattform stattgefunden. Es wurde
    in der G20-Plattform eine konkrete Arbeitsweise entwickelt.

    Es wird eine Demo im April geben. Bis dahin soll mobilisiert werden.

    Der AStA informiert sich über Chancen eines Volksentscheids. Am 24.1. findet
    um 19 Uhr das nächste Treffen statt.

    Artur verweist auf eine Veranstaltung namens "Putsch im Putsch in der Türkei
    und Berufsverbote in Deutschland" am 21. Januar 19 Uhr.

    Geoffrey setzt zum Ende des Berichts an:
    Das neue Stille Wörtchen ist fertig mit dem Schwerpunkt Wahlen. Es enthält
    einen Rückblick, Infos zum Akademischen Senat und weitere Infos dazu, was so
    gewählt wird. Anwesende von Außenstandorten sind eingeladen das StiWö
    mitzunehmen und aufzuhängen.

    \subsubsection{Anfragen an den AStA}

    Tobias B. fragt Armin, wie die Gesellschaft überzeugt werden soll.
    Welche Bündnispartner sind angedacht? Wie wird mit DGB-Entscheidung umgegangen?

    Armin antwortet:
    Manifest ist ein guter Start, um es in die Gesellschaft zu tragen.
    AStA-Infobroschüre zum G20 kann auch hilfreich dahingehend sein. Es sei ein
    langer Weg bis zum gesamtgesellschaftlichem Widerstand. Man müsse Mitglieder
    im DGB und in Parteien motivieren gegen Vorstandsentscheidung anzugehen.

    Franziska ergänzt:
    Es wird in irgendeiner Form eine Unterschriftenkampagne geben.

    \subsubsection{Dringlichkeitsanträge des AStA}

    Keine.

    \newpage
    \subsubsection{Feststellung der endgültigen Fassung des Teils B der Tagesordnung}

    %%% BESCHLOSSENE TAGESORDNUNG %%%
    \underline{Endgültige Fassung der Tagesordnung}

    Es wird der Antrag gestellt, das Manifest an erster Stelle zu behandeln.
    Dieser Antrag wird abgelehnt.

    Ein weiterer Antrag möchte das Manifest an zweiter Stelle behandeln. Dies
    wird angenommen.
    \begin{enumerate}[label={\textbf{Top \theenumi}},leftmargin=*]
        \item Beitragsordnung (V1617-046) (30 Min.)
        \item Manifest zu G20 (V1617-049) (45 Min.)
        \item VG-Wort (V1617-047) (30 Min.)
        \item Englischsprachige Wahlwerbung (V1617-048) (30 Min.)
        \item Verschiedenes (5 Min.)
    \end{enumerate}

    Die so geänderte Tagesordnung wird mehrheitlich bei 0 Gegenstimmen und 2
    Enthaltungen angenommen.

    \subsubsection{Feststellung der Beschlussfähigkeit}

    \textit{Mit 28 Parlamentarier*innen beschlussfähig.}

    \subsubsection{Genehmigung der Protokolle der vorangegangenen Sitzungen}

    Keine Protolle zu genehmigen. Die ausstehenden Protokolle sollen auf der
    14. Sitzu

    \vspace{0.5cm}
    {\Large \underline{Teil B}}

    \textit{Zeitansage: 21.43 Uhr}

    \subsection{Beitragsordnung}

    Philipp stellt den Antrag vor.

    Franziska stellt vor, dass eine 0,10€ Verringerung des AStA-Budgets eine
    Veränderung von ca. 8400€ bedeutet.

    Tobias B. fragt, ob der AStA sich überlegt hat, die Forderungen des HVV nicht
    zu bedienen. Philipp sagt, dass der AStA sich frühzeitiger engagiert hat,
    dies aber nicht gefruchtet hat.

    Gunhild meint, dass frühere ASten offensiver in der Frage agierten.

    Artur merkt an, dass nach längerer Diskussion die Vertreter des HVV die
    Punkte des AStA im HVV intern diskutieren wollten. Bei einer "Busbeschleunigungsfahrt"
    wurde mit dem zuständigen Staatsrat gesprochen und der hat sich gegen die Verteilung
    des Senatsprozents an die Studierenden ausgesprochen, weil dann ja alle kommen
    könnten.

    Jennifer bittet, dass sich in der Debatte um Beitragsordnung nicht über
    kostenlosen ÖPNV unterhalten wird.

    Tobias B. findet, dass man sich zusammen mit den anderen ASten ab sofort
    zusammenschließen sollte, um eine Mobilisierung gegen die Erhöhungen zu
    erreichen. Tobias B. fragt, warum in der Vergangenheit nichts dahingegend
    passiert ist.

    Jennifer fragt, warum denn kein entsprechender Antrag gestellt wurde.

    Till erinnert an aktivere Zeiten. Es muss durch Öffentlichkeitsarbeit und
    Skandalisierung kontinuierlich auf das Problem hingewiesen wird. Die Botschaft
    sollte lauten: Semesterticket ja, aber nicht um jeden Preis.
    In der Vergangenheit gab es dazu Anträge und Beschlüsse. Die sollte man
    wieder herausholen.

    Gunhild merkt an, dass es mal einen Ausschuss zu dem Semesterticket gab.

    Philipp sagt, dass es den Versuch von allen Hamburger ASten gab am Ball zu
    bleiben, dies aber auf der Strecke geblieben ist. Es gibt bereits Arbeit,
    an der angedockt werden kann.

    Artur merkt an, dass ein Begründungsteil erreicht wurde in dem Beitragsbescheid,
    der Infos zum Senatsprozent enthält.

    Der Antrag wird mehrheitlich bei einer Gegenstimme und keiner Enthaltung
    angenommen.

    \subsection{Manifest zu G20}

    Artur stellt den Antrag vor.

    Yasemin fragt, warum gefordert wird die Tarife und Personalmittel im
    öffentlichen Dienst zu erhöhen, wo es durchaus Menschen gibt, die weit mehr
    als Mindestlohn verdienen.

    Tobias B. entgegnet, dass eine Erhöhung dort allen Beschäftigten zugute komme
    und ein Verteilungskampf nicht gestartet werden sollte.

    \textit{Zeitansage: 22.11 Uhr}

    Franziska merkt an, dass Bezahlung im öffentlichen Dienst eher unter dem Einfluss
    von Aktiven steht als privatwirtschaftliche Verhältnisse. Es kann als Vorbild
    fungieren.

    Artur weist darauf hin, dass Personalmittel auch für mehr Einstellungen
    stehen und dies sinnvoll ist, weil diese Gruppen in der Vergangenheit eine
    große Arbeitsverdichtung erfahren haben.

    Tobias B. fragt, warum Yasemin die geäußerte Position hat.

    Tobias B. merkt an, dass die Stadt der größte Arbeitgeber in der Stadt ist.
    Die Bundesbank fordert mittlerweile bereits, dass Gewerkschaften echte
    Lohnforderungen stellen sollen. Ebenfalls weist er darauf hin, dass Deutschland
    ein Exportüberschuss von rund 200 Mrd. Euro hat und die in Lohnforderungen
    umgesetzt werden sollten.

    Timo hat nachgerechnet, dass für die Elbphilharmonie allen Kindern 2,5 Jahre
    Musikunterricht hätte gegeben werden können. Diese Menge Geld sollte daher
    jetzt nicht auch noch in den G20-Gipfel gesteckt werden.

    Franziska meint, dass ein Großteil der Bevölkerung Ahnung habe, wo Geld
    investiert werden könne.

    Tobias B. stellt die Frage in den Raum, wie das Staatswesen verfasst sein
    sollte. Prestige oder echte Arbeit an Verbesserung? Man müsste verstehen, dass
    der Topf Personalmittel einer Stadt wichtig ist, da wichtige Dinge Geld
    kosten und die Arbeitenden ihre Aufmerksamkeit auf sinnvolle Arbeit lenken
    können sollten und nicht auf ihren eigenen Haushalt.
    Tobias erwähnt ferner, dass es wert ist das Geld an anderer Stelle wegzunehmen.

    Franziska erwähnt, dass alle aufgerufen sind mit dem Manifest weitere Unterstützer
    zu suchen.

    Tobias B. fragt, wie man sich auf das Manifest beziehen kann.

    Franziska antwortet, dass das Manifest für Demomobilisierung verwendet werden
    soll. Eine Pressekonferenz wird es dann im März geben.

    Der Antrag wird mehrheitlich bei einer Gegenstimme und 3 Enthaltungen
    angenommen.

    \textit{Zeitansage: 22.29 Uhr}
    \subsection{VG-Wort}

    Geoffrey stellt den Antrag vor.

    Till wiederholt seine Anmerkungen von der 12. Sitzung. Er hat weiterhin Probleme
    mit dem Antrag und liest seine Alternativversion vor. Till schlägt vor
    für die nächste Sitzung einen Konsensvorschlag einzureichen.

    Karima findet einen Angriff auf Digitalisierung problematisch.

    Till findet, dass eine Forderung fehlt, dass in Bibliotheken genügend
    Materialien vorhanden sein sollen.

    \textit{Der TOP wird auf die nächste Sitzung vertagt.}

    \subsection{Englischsprachige Wahlwerbung}

    Jim stellt den Antrag vor.

    In der entfachten und teils hitzigen Debatte werden im Grunde zwei Positionen
    vertreten. Die eine tritt für englischsprachige Wahlwerbung ein, während die
    andere scheinbar dagegen ist. Bei näherem Hinsehen fällt jedoch auf, dass
    beide für solche Werbung sind, die zweite Position jedoch Realisierungsprobleme
    sieht und vor allem daraufhin orientieren möchte, dass die betroffenen
    Studierenden bessere Deutsch als Fremdsprache Angebote bekommen.

    Über diese inhaltliche Diskussion wird jedoch auch eine Metadiskussion
    entzündet, welche den schwelenden Konflikt zwischen Liste Links und harte Zeiten
    auf der einen und der Alternativen Linken auf der anderen Seite offenbart.
    Die Gesprächskultur leidet darunter.

    Timo Tomanek macht einen Vorschlag zur Güte und fragt sich, warum nicht
    sowohl englischsprachige Wahlwerbung gemacht und sich für bessere
    Angebote von Deutsch als Fremdsprache eingesetzt werden kann.

    \textit{Melf stellt GO-Antrag zum Ende der Debatte. Als Alternative wird
    vorgeschlagen die Redeliste zu schließen. Es gibt 8 Stimmen für das Ende
    der Debatte, 10 dagegen und 2 Personen enthalten sich. Somit wird die
    Redeliste geschlossen und abgearbeitet.}

    \textit{Geoffrey stellt den GO-Antrag die Sitzung nach dem TOP abzubrechen.
    Dieser wird mehrheitlich ohne Gegenstimmen bei einer Enthaltung angenommen.}

    \textit{Die Sitzung wird um 23.42 Uhr abgebrochen}.
\end{document}
